\documentclass[10pt]{article}
\usepackage[a4paper,margin=0.6in]{geometry}
\usepackage{amsmath,amssymb,amsfonts,mathtools}
\setlength{\jot}{4pt}
\usepackage{xcolor}
\setlength{\parindent}{0pt}
\setlength{\parskip}{5pt}
\newcommand{\Leps}{L_{\epsilon}}
\title{Group II, Row IV: reduction of the master expression\\ to the three-term ($\mathcal P_{gg}$) form}
\author{}
\date{}
\begin{document}
\maketitle
\vspace{-2.5em}

\textbf{Claim.} The master expression
\begin{align}\label{eq:master}
     \frac{d^{3}N_{\rm NLO}^{\rm}}{d^{3}k}  & = \frac{1}{(2\pi)^{3}}\,\frac{g^{4}}{4\pi^{4}}\,\int_{\Lambda}^{\vee-k^{+}}\, \frac{dp^{+}}{2\pi}\int_{z,w,w^{\prime},x'} e^{-ik(w^{\prime}-w)}\,  \int_{x,y,y^{\prime}}\,\frac{(w-z)^{m}\,(x-y)^{k}\,(w'-z)^{m'}\,(x'-y')^{k'}}{(w-z)^{2}\,(x-y)^{2}\,(w'-z)^{2}\,(x'-y')^{2}}\,\frac{p^{+}\,k^{+}}{(p^{+}+k^{+})^{4}}\notag \\ & \delta^{(2)}\left[y-w+\frac{p^{+}}{(p^{+}+k^{+})}\,(w-z)\right]\, \delta^{(2)}\left[y'-w'+\frac{p^{+}}{(p^{+}+k^{+})}\,(w'-z)\right]\times\notag \\ &  \left[\delta_{km}\,\delta_{ij}-\frac{(p^{+}+k^{+})}{p^{+}}\,\delta_{jm}\,\delta_{ik}-\frac{(p^{+}+k^{+})}{k^{+}}\,\delta_{im}\,\delta_{kj}\right] \left[\delta_{k'm'}\,\delta_{ij}-\frac{(p^{+}+k^{+})}{p^{+}}\,\delta_{jm'}\,\delta_{ik'}-\frac{(p^{+}+k^{+})}{k^{+}}\,\delta_{im'}\,\delta_{k'j}\right]\notag \\ & \left[f^{c'd'a}\,U^{c'e'}(y')\,U^{d'b}(z)\,\rho^{e'}(x')-f^{c'd'a}\, U^{d'b}(z)\,U^{c'e'}(x')\,\rho^{e'}(x')\right] \left[f^{cda}\,U^{ce}(y)\,U^{db}(z)\,\rho^{e}(x)-f^{cda}\, U^{db}(z)\,U^{ce}(x)\,\rho^{e}(x)\right]
\end{align}
is identically equal, up to $O(\Lambda)$ and $O(\epsilon)\times{\rm finite}$, to the three-term form
\begin{align}\label{eq:target}
    \frac{d^{3}N_{\rm NLO}^{\rm}}{d^{3}k}  & = \frac{g^{2}}{4\pi^{2}}\, \frac{N_{c}}{(2\pi)^{3}}\,\frac{g^{2}}{2\pi^{2}}\, \frac{1}{k^{+}}\,\int_{x,y,x',y'}\,\frac{(x-y)\cdot (x'-y')}{(x-y)^{2}\,(x'-y')^{2}}\,\left[-\frac{d_{\perp}}{\epsilon}-d_{\perp}\,\gamma_{E}-\frac{d_{\perp}}{2}\,\log{\left(\frac{(y'-y)^{2}\,\mu^{2}_{\overline{MS}}}{4}\right)}-1\right]\,\notag \\ & \Bigg[\int_{\frac{k^{+}}{\vee}}^{1}\,\frac{d\xi}{\xi^{2}}\,e^{-ik\cdot (y'-y)\,\frac{1}{\xi}}\,\left[\xi\,(1-\xi)+\frac{(1-\xi)}{\xi}+\frac{\xi}{(1-\xi)_{+}}\right]\Bigg]\times \left[U^{ce'}(y') - U^{ce'}(x')\right] \left[U^{ce}(y) - U^{ce}(x)\right] \rho^{e'}(x') \rho^{e}(x)\notag \\ &  +\, \frac{g^{2}}{4\pi^{2}}\, \frac{1}{(2\pi)^{3}}\,\frac{g^{2}}{2\pi^{2}}\, \frac{N_{c}}{k^{+}}\,{\color{blue}\left[\log{\frac{(\vee - k^{+})\, k^{+}}{\vee\,\Lambda}}\right]}\,\int_{x,y,x',y'}\,e^{-ik\cdot (y'-y)}\,\frac{(x-y)\cdot (x'-y')}{(x-y)^{2}\,(x'-y')^{2}}\,\bigg[-\frac{d_{\perp}}{\epsilon}-d_{\perp}\,\gamma_{E}\notag \\ & -\frac{d_{\perp}}{2}\,\log{\left(\frac{(y'-y)^{2}\,\mu^{2}_{\overline{MS}}}{4}\right)}-1\bigg]\,\left[U^{ce'}(y') - U^{ce'}(x')\right] \left[U^{ce}(y) - U^{ce}(x)\right] \rho^{e'}(x') \rho^{e}(x)\notag \\ & +\,\frac{g^{2}}{4\pi^{2}}\, \frac{1}{(2\pi)^{3}}\,\frac{g^{2}}{2\pi^{2}}\, \frac{N_{c}}{k^{+}}\,\int_{\frac{k^{+}}{\vee}}^{\frac{k^{+}}{k^{+}+\Lambda}}\,d\xi\, \frac{(1-\xi)}{\xi}\, \int_{x,y,x',y'}\,e^{-ik\cdot (y'-y)\frac{1}{\xi}}\,\Bigg[\frac{(x-y)\cdot (x'-y')}{(x-y)^{2}\, (x'-y')^{2}} -\frac{(x-y)\cdot (y'-y)\, (x'-y')\cdot (y'-y)}{(x-y)^{2}\,(x'-y')^{2}\,(y'-y)^{2}}\,d_{\perp}\Bigg]\notag \\ &  \times\left[U^{ce'}(y') - U^{ce'}(x')\right] \left[U^{ce}(y) - U^{ce}(x)\right] \rho^{e'}(x') \rho^{e}(x)\,,
\end{align}
where in the first term $\frac{\xi}{(1-\xi)_{+}}$ denotes the \emph{locally subtracted} prescription,
$\int_{k^{+}/\vee}^{1}d\xi\,\frac{h(\xi)}{(1-\xi)_{+}}\equiv\int_{k^{+}/\vee}^{1}d\xi\,\frac{h(\xi)-h(1)}{1-\xi}$,
whose extracted endpoint is the second (blue) term.

\section{Colour algebra}
The two colour brackets factor as $f^{c'd'a}U^{d'b}(z)[U^{c'e'}(y')-U^{c'e'}(x')]\rho^{e'}(x')$ and
$f^{cda}U^{db}(z)[U^{ce}(y)-U^{ce}(x)]\rho^{e}(x)$. Using the reality/orthogonality of the adjoint
Wilson line, $U^{d'b}(z)U^{db}(z)=\delta^{d'd}$, and then $f^{c'da}f^{cda}=N_{c}\,\delta^{c'c}$,
\begin{align}
(\text{colour brackets})\;=\;N_{c}\,\big[U^{ce'}(y')-U^{ce'}(x')\big]\big[U^{ce}(y)-U^{ce}(x)\big]\,\rho^{e'}(x')\rho^{e}(x)\,.
\end{align}

\section{$\delta$-function integration over $w,w'$}
With $\lambda=\frac{p^{+}}{p^{+}+k^{+}}$ and $1-\lambda=\frac{k^{+}}{p^{+}+k^{+}}$,
\begin{align}
\delta^{(2)}\big[y-(1-\lambda)w-\lambda z\big]=\Big(\frac{p^{+}+k^{+}}{k^{+}}\Big)^{2}\,
\delta^{(2)}\Big[w-\frac{(p^{+}+k^{+})y-p^{+}z}{k^{+}}\Big]\,,
\end{align}
and at the solution
\begin{align}
w-z=\frac{p^{+}+k^{+}}{k^{+}}\,(y-z)\,,\qquad w'-w=\frac{p^{+}+k^{+}}{k^{+}}\,(y'-y)\,.
\end{align}
Hence, per side, $\int d^{2}w\,\delta^{(2)}\,\frac{(w-z)^{m}}{(w-z)^{2}}
=\frac{p^{+}+k^{+}}{k^{+}}\,\frac{(y-z)^{m}}{(y-z)^{2}}$, the phase becomes
$e^{-ik(y'-y)\frac{p^{+}+k^{+}}{k^{+}}}$, and the prefactor collects to
\begin{align}
\frac{p^{+}k^{+}}{(p^{+}+k^{+})^{4}}\,\Big(\frac{p^{+}+k^{+}}{k^{+}}\Big)^{2}
=\frac{p^{+}}{k^{+}\,(p^{+}+k^{+})^{2}}\,.
\end{align}

\section{Tensor contraction over $i,j$}
With $A=\frac{p^{+}+k^{+}}{p^{+}}$, $B=\frac{p^{+}+k^{+}}{k^{+}}$, a direct index computation
(checked mechanically over all index values) gives
\begin{align}
\sum_{i,j}T_{1}^{(km,ij)}T_{2}^{(k'm',ij)}
=\big[d_{\perp}-2(A+B)\big]\,\delta_{km}\delta_{k'm'}+(A^{2}+B^{2})\,\delta_{mm'}\delta_{kk'}
+2AB\,\delta_{mk'}\delta_{km'}\,,
\end{align}
and multiplying by $\frac{p^{+}}{k^{+}(p^{+}+k^{+})^{2}}$,
\begin{align}
-2(A+B)\to-\frac{2}{k^{+2}}\,,\qquad
A^{2}+B^{2}\to\frac{1}{k^{+}}\Big(\frac{1}{p^{+}}+\frac{p^{+}}{k^{+2}}\Big)\,,\qquad
2AB\to+\frac{2}{k^{+2}}\,,\qquad
d_{\perp}\to\frac{d_{\perp}\,p^{+}}{k^{+}(p^{+}+k^{+})^{2}}\,.
\end{align}
Since the $z$-integral below is a \emph{symmetric} tensor in $(m,m')$, the
$\delta_{km}\delta_{k'm'}$ and $\delta_{mk'}\delta_{km'}$ channels contract identically with the
kernels, so \textbf{the $\mp\frac{2}{k^{+2}}$ pieces cancel exactly}: no antisymmetric channel
survives.

\section{The $z$ integral}
Only $\mathcal K^{m}(y-z)\mathcal K^{m'}(y'-z)$ depends on $z$. In $d_{\perp}=2-\epsilon$
($\overline{\rm MS}$), with $r=y'-y$,
\begin{align}\label{eq:zint}
\int_{z}\,\frac{(y-z)^{m}\,(y'-z)^{m'}}{(y-z)^{2}\,(y'-z)^{2}}
=-\frac{\pi}{2}\,\underbrace{\Big[\frac{2}{\epsilon}+2\gamma_{E}+\log\frac{r^{2}\mu^{2}_{\overline{MS}}}{4}\Big]}_{\textstyle \Leps}\,\delta^{mm'}\,-\,\pi\,\frac{r^{m}r^{m'}}{r^{2}}\,.
\end{align}
\emph{Independent check:} the traceless projection of the left-hand side is finite in $d=2$ and was
evaluated numerically: it equals $-\pi\big[\frac{r^{m}r^{m'}}{r^{2}}-\frac{\delta^{mm'}}{2}\big]$ to
$10^{-6}$; the trace grows with the infrared cutoff with the $r$-dependence $-\pi\log r^{2}$, fixing
the structure of \eqref{eq:zint} (the constant defines $\mu_{\overline{MS}}$). Contracting with the
surviving channels of Section~3 and with $\mathcal K^{k}(x-y)\mathcal K^{k'}(x'-y')$:
\begin{itemize}
\item channel $\frac{d_{\perp}p^{+}}{k^{+}(p^{+}+k^{+})^{2}}\,\delta_{km}\delta_{k'm'}$:
$\;-\frac{\pi}{2}\Leps\,\mathcal K(x-y)\!\cdot\!\mathcal K(x'-y')
-\pi\,\frac{[\mathcal K(x-y)\cdot r][\mathcal K(x'-y')\cdot r]}{r^{2}}\,$ (the azimuthal kernel, carrying $d_{\perp}$);
\item channel $\frac{1}{k^{+}}\big(\frac{1}{p^{+}}+\frac{p^{+}}{k^{+2}}\big)\,\delta_{mm'}\delta_{kk'}$:
trace of \eqref{eq:zint} times the dipole,
$\;\big[-\frac{\pi}{2}d_{\perp}\Leps-\pi\big]\,\mathcal K(x-y)\!\cdot\!\mathcal K(x'-y')$,
where the extra $-\pi$ comes from $\delta_{mm'}\frac{r^{m}r^{m'}}{r^{2}}=1$.
\textbf{This $-\pi$ is the origin of the ``$-1$'' in the bracket of \eqref{eq:target}.}
\end{itemize}

\section{$\xi$ substitution and matching}
Set $\xi=\frac{k^{+}}{p^{+}+k^{+}}$, so that
\begin{align}
p^{+}=\frac{k^{+}(1-\xi)}{\xi},\quad dp^{+}=\frac{k^{+}d\xi}{\xi^{2}},\quad
e^{-ik(y'-y)\frac{p^{+}+k^{+}}{k^{+}}}=e^{-ik(y'-y)/\xi},\quad
p^{+}\in(\Lambda,\vee-k^{+})\leftrightarrow\xi\in\Big(\frac{k^{+}}{\vee},\frac{k^{+}}{k^{+}+\Lambda}\Big),
\end{align}
\begin{align}
\frac{p^{+}}{(p^{+}+k^{+})^{2}}=\frac{\xi(1-\xi)}{k^{+}},\qquad
\frac{p^{+}}{k^{+2}}=\frac{1-\xi}{\xi k^{+}},\qquad
\frac{1}{p^{+}}=\frac{\xi}{(1-\xi)k^{+}}\,.
\end{align}
The $\Leps$ coefficient assembles into
$\frac{d_{\perp}}{2}\,\hat p(\xi)$ with
$\hat p(\xi)=\xi(1-\xi)+\frac{1-\xi}{\xi}+\frac{\xi}{1-\xi}$, and the overall constant is
$\frac{1}{(2\pi)^{3}}\frac{g^{4}N_{c}}{8\pi^{4}}\frac{1}{k^{+}}
=\frac{g^{2}}{4\pi^{2}}\frac{N_{c}}{(2\pi)^{3}}\frac{g^{2}}{2\pi^{2}}\frac{1}{k^{+}}$, as in
\eqref{eq:target}. Altogether,
\begin{align}\label{eq:intermediate}
\frac{d^{3}N_{\rm NLO}}{d^{3}k}
=-\,C\!\!\int_{\frac{k^{+}}{\vee}}^{\frac{k^{+}}{k^{+}+\Lambda}}\!\!\!d\xi\,
\Big\{\frac{d_{\perp}}{2}\,\frac{\hat p(\xi)}{\xi^{2}}\,\Leps
+\frac{1}{\xi^{2}}\Big[\frac{\xi}{1-\xi}+\frac{1-\xi}{\xi}\Big]\Big\}\,
\mathcal K\!\cdot\!\mathcal K'\;f(\xi)
\;-\;C\,d_{\perp}\!\!\int\! d\xi\,\frac{1-\xi}{\xi}\,
\frac{(\mathcal K\!\cdot\! r)(\mathcal K'\!\cdot\! r)}{r^{2}}\,f(\xi)\,,
\end{align}
with $C=\frac{g^{2}}{4\pi^{2}}\frac{N_{c}}{(2\pi)^{3}}\frac{g^{2}}{2\pi^{2}}\frac{1}{k^{+}}$,
$f(\xi)=e^{-ik(y'-y)/\xi}$, $\mathcal K=\mathcal K(x-y)$, $\mathcal K'=\mathcal K(x'-y')$, the colour
brackets understood. Two exact regroupings turn \eqref{eq:intermediate} into \eqref{eq:target}:

\textbf{(a) The endpoint.} For the $\frac{\xi}{1-\xi}$ channel, with $h(\xi)$ the rest of the integrand,
\begin{align}
\int_{\frac{k^{+}}{\vee}}^{\frac{k^{+}}{k^{+}+\Lambda}}\!\!d\xi\,\frac{h(\xi)}{1-\xi}
=\int_{\frac{k^{+}}{\vee}}^{1}\!\!d\xi\,\frac{h(\xi)-h(1)}{1-\xi}
\;+\;h(1)\,\log\frac{(\vee-k^{+})\,k^{+}}{\vee\,\Lambda}\;+\;O(\Lambda)\,,
\end{align}
which is exactly the locally subtracted $\frac{\xi}{(1-\xi)_{+}}$ of the first term of
\eqref{eq:target} plus the blue endpoint term, both carrying the same bracket
$-\big[\frac{d_{\perp}}{2}\Leps+1\big]$.

\textbf{(b) The ``$-1$'' and the azimuthal $\mathcal K\!\cdot\!\mathcal K'$ piece.} The target writes the
$\Leps$-free dipole piece as ``$-1$'' against the \emph{full} $\hat p/\xi^{2}$ plus the
$+\mathcal K\!\cdot\!\mathcal K'\,\frac{1-\xi}{\xi}$ term inside the third line of \eqref{eq:target}.
This is identical to the $\Leps$-free piece of \eqref{eq:intermediate} because of the exact identity
\begin{align}
\frac{\hat p(\xi)}{\xi^{2}}\;-\;\frac{1-\xi}{\xi}
\;=\;\frac{1}{\xi^{2}}\Big[\frac{\xi}{1-\xi}+\frac{1-\xi}{\xi}\Big]\,.
\end{align}
Finally, the $r$-tensor part of \eqref{eq:intermediate} carries exactly
$-\,d_{\perp}\frac{1-\xi}{\xi}$, matching the second kernel of the third term of \eqref{eq:target}.

\section*{Conclusion}
Equations \eqref{eq:master} and \eqref{eq:target} are identical, up to $O(\Lambda)$ (from the endpoint
recombination and from extending regular $\xi$-integrals to $\xi=1$) and
$O(\epsilon)\times\text{finite}$. All coefficients check, including: the overall $N_{c}$; the
``$-1$'' in the dimensional-regularization bracket (the trace of the $r$-tensor part of the
$z$-integral against the $\frac{1}{p^{+}}+\frac{p^{+}}{k^{+2}}$ channel); the $d_{\perp}$ on the
azimuthal kernel; and the complete cancellation of the antisymmetric
$\frac{2}{k^{+}}(\delta_{mk'}\delta_{m'k}-\delta_{m'k'}\delta_{km})$ channel between the two crossed
contractions (the $z$-integral is symmetric in $m,m'$).

\emph{Convention remark.} The index algebra is carried out in $d_{\perp}=2-\epsilon$ while the
position-space measures and the $\delta$-functions are kept two-dimensional; only the $z$ loop is
dimensionally continued. This is consistent as long as it is applied uniformly --- as it is on both
sides of the equality --- since any $\xi^{2\epsilon}$ Jacobian refinement would only shift
scheme-dependent finite pieces that cancel in the $\mu$-independent total.
\end{document}
